% Einheit 3 von 4 – Einheit 3 (bank/lineare-gleichungen/e3.jsonl, zusammenbau v0.1)
%% TODO Prüfungswort Einheit 3: nicht in der Marken-Zeile
%% TODO Zeitmarke Einheit 3: Marken-Zeile nicht lesbar
%% TODO Zweigzeile Teil 1: was hier gelernt wird (Satz in Schülersprache aus dem Einheitstitel) – Einheit 3
%% TODO Einheitentitel 3 nicht in der Mappe lesbar
\einheitenkopf[e3]{Einheit 3 von 4 · Einheit 3}
\setcounter{aufgabe}{15}

%% TODO Titel: Ich-kann-Satz fehlt in der Bank (Platzhalter: Kettenname) – Nr. 16
%% TODO Anweisung: ein Satz über der ersten Teilaufgabe fehlt in der Bank – Nr. 16
\begin{aufgabe}{x beidseitig}
\begin{teile}
\teil $7x + 2 = 3x + 14$ – wo steht weniger $x$? \leerfeld
\end{teile}
%% TODO Darstellung für schwach fehlt in der Bank (lineare-gleichungen-e3-k1-s1-v1; Abschnitt „Für schwache Schüler“)
\swz{$5x = 2x + 12$}{}
%% TODO Darstellung für schwach fehlt in der Bank (lineare-gleichungen-e3-k1-s1-v2; Abschnitt „Für schwache Schüler“)
\swz{$7x = 3x + 20$}{}
%% TODO Darstellung für schwach fehlt in der Bank (lineare-gleichungen-e3-k1-s1-v3; Abschnitt „Für schwache Schüler“)
\swz{$6x = x + 30$}{}
%% TODO Darstellung für schwach fehlt in der Bank (lineare-gleichungen-e3-k1-s1-v4; Abschnitt „Für schwache Schüler“)
\swz{$9x = 4x + 35$}{}
%% TODO Darstellung für schwach fehlt in der Bank (lineare-gleichungen-e3-k1-s1-v5; Abschnitt „Für schwache Schüler“)
\swz{$4x = x + 24$}{}
\swfrage{Was bleibt gleich, was ändert sich?}
%% TODO Darstellung für schwach fehlt in der Bank (lineare-gleichungen-e3-k1-s2-v1; Abschnitt „Für schwache Schüler“)
\swz{$6x + 5 = 2x + 25$}{}
%% TODO Darstellung für schwach fehlt in der Bank (lineare-gleichungen-e3-k1-s3-v1; Abschnitt „Für schwache Schüler“)
\swz{$4x + 7 + x = 2x + 22$}{}
%% TODO Darstellung für schwach fehlt in der Bank (lineare-gleichungen-e3-k1-s4-v1; Abschnitt „Für schwache Schüler“)
\swz{$3(x + 2) = x + 14$}{}
%% TODO Darstellung für schwach fehlt in der Bank (lineare-gleichungen-e3-k1-s5-v1; Abschnitt „Für schwache Schüler“)
\swz{$\frac{x}{2} + 4 = x - 1$}{}
%% TODO Darstellung für schwach fehlt in der Bank (lineare-gleichungen-e3-k1-s6-v1; Abschnitt „Für schwache Schüler“)
\swz{$2(x + 3) = 2x + 6$}{}
\swz[3]{$\text{$5(x - 2) = 3x + 4$ – löse und mache die Probe.}$}{}
\end{aufgabe}

%% TODO Titel: Ich-kann-Satz fehlt in der Bank (Platzhalter: Kettenname) – Nr. 17
%% TODO Anweisung: ein Satz über der ersten Teilaufgabe fehlt in der Bank – Nr. 17
\begin{aufgabe}{Dezimalzahlen}
%% TODO Darstellung für schwach fehlt in der Bank (lineare-gleichungen-e3-k2-s1-v1; Abschnitt „Für schwache Schüler“)
\swz{$0{,}5x + 1{,}2 = 3{,}7$}{}
\end{aufgabe}

%% TODO Titel: Ich-kann-Satz fehlt in der Bank (Platzhalter: Kettenname) – Nr. 18
%% TODO Anweisung: ein Satz über der ersten Teilaufgabe fehlt in der Bank – Nr. 18
\begin{aufgabe}{x beidseitig – Fehler finden, Begründen}
\swz[3]{Nina rechnet: \rechnung{2x + 5 + 4x &= 23 \\ 11x &= 23} Finde den Fehler.}{}
\swz[3]{Ole rechnet: \rechnung{15 - (x - 3) &= 2x \\ 15 - x - 3 &= 2x \\ 12 &= 3x \\ 4 &= x} Finde den Fehler.}{}
\swz[3]{Paul rechnet: \rechnung{3(x + 1) &= 3x + 3 \\ 3x + 3 &= 3x + 3 &&\mid -3x \\ 3 &= 3} Er schreibt: „Die Lösung ist 0." Finde den Fehler.}{}
\swfrage{$x + 4 = x + 9$ – gibt es eine Lösung? Begründe ohne Rechnung.}
\end{aufgabe}

\uebersichtskasten{x auf beiden Seiten: erst alle x auf eine Seite (die mit mehr x), dann Zahlen auf die andere. \\ 5x + 2 = 3x + 10 | −3x → 2x + 2 = 10 \\ Klammern zuerst auflösen, dann zusammenfassen, dann umformen. \\ Brüche: mit dem Hauptnenner malnehmen. \quad x/3 + 2 = 4 | ·3 → x + 6 = 12 \\ Sonderfälle: 0 = 5 → keine Lösung \quad 0 = 0 → jede Zahl ist Lösung}
